\section{Normalized spherical coefficients on common elliptic torsors}
\label{sec:spherical-hall}

There is a smooth sector in which the relative coefficients and their
extension maps can be constructed explicitly. Its Hall product retains
the diagonal of the elliptic curve. Direct sum gives a coproduct on the
same coefficient objects, but compatibility with multiplication requires
a symmetry that also retains this diagonal. We first construct the
objects and the two geometric operations. Their global comparison
is Theorem~\ref{gh:global-hopf}; Section~\ref{sec:supported-springer}
constructs the crossing itself over support from proper full flags.

\subsection{The smooth spherical sector}

Fix a spherical coherent sheaf $G$ on the projective K3 surface $S$:
\[
 \operatorname{Hom}_S(G,G)=\mathbb C,\qquad
 \operatorname{Ext}^1_S(G,G)=0,\qquad
 \operatorname{Ext}^2_S(G,G)=\mathbb C.
\]
Choose a holomorphic volume form on $S$ and a nonzero differential on
$E$; these fix the Serre traces. The example $G=\mathcal O_S$ is always
available. Let $\mathcal M_n=\mathfrak T_n(E)$ be the stack of
length-$n$ torsion sheaves on $E$, and put
\[
 B_E=BE,\qquad
 \overline{\mathcal M}_n=[\mathcal M_n/E],\qquad
 \pi_n:\overline{\mathcal M}_n\longrightarrow B_E.
\]
The charge is the length $n$. In particular
$\mathcal M_0=\operatorname{Spec}\mathbb C$ and
$\overline{\mathcal M}_0=B_E$.

\begin{proposition}\label{gh:sector}
The functor $T\mapsto G\boxtimes T$ is exact and fully faithful, and its
image is closed under extensions in $\operatorname{Coh}(S\times E)$.
It identifies $\mathcal M_n$ with a smooth open sheaf locus of dimension
zero. For every ordered list $\mathbf n=(n_1,\ldots,n_r)$, the full flag
stack $\mathcal F_{\mathbf n}$ has a smooth endpoint map of relative
dimension zero and a representable projective map to
$\mathcal M_{\sum n_i}$. After simultaneous translation quotient, the
flag and middle-object stacks have dimension $-1$. These assertions
include nonsplit flags and zero blocks.
\end{proposition}
\begin{proof}
K\"unneth gives, for torsion sheaves $T,U$,
\[
 R\operatorname{Hom}_{S\times E}(G\boxtimes T,G\boxtimes U)
 \simeq R\operatorname{Hom}_S(G,G)
             \otimes R\operatorname{Hom}_E(T,U).
\]
In degrees zero and one this identifies Hom and Ext with those on $E$.
Thus the functor is fully faithful and every ambient extension class
comes from an extension on $E$. The two middle terms are then
isomorphic. The same comparison holds in families.

For completeness, a length-$n$ sheaf is a quotient of
$\mathcal O_E^n$ whose map on global sections is an isomorphism.
Let $Q_n$ be this open part of the Quot scheme. Its kernel is a vector
bundle $K$ of rank $n$, and
$\operatorname{Ext}^1(K,T)=H^1(K^\vee\otimes T)=0$.
The tangent space $\operatorname{Hom}(K,T)$ has dimension $n^2$.
Hence $\mathcal M_n=[Q_n/\operatorname{GL}_n]$ is smooth of dimension
zero. The exterior-product map is a monomorphism of sheaf stacks.
Across a square-zero extension, the source is unobstructed and the
Hom and Ext$^1$ comparisons identify automorphisms and possible lifts.
Induction on small extensions therefore makes this map formally
\'etale. Being locally of finite presentation, it is an open immersion.

The groupoid of extensions of fixed endpoints is locally the vector
bundle stack $[V^1/V^0]$ associated with
$R\operatorname{Hom}_E(T_b,T_a)[1]$.
There is no Ext$^2$ on a smooth curve, and the Euler pairing of two
torsion sheaves is zero, so the two vector bundles have equal rank.
This proves smoothness and relative dimension zero. Iteration gives the
full endpoint map. The map forgetting a subobject is the relative Quot
map and is projective and representable. Equivalently, on a finite
support presentation it is the closed condition that a subspace in a
Grassmannian be stable under the support algebra. This description
includes nilpotent support. Iterating it proves the full flag assertion.
Quotienting all terms by the same smooth one-dimensional group subtracts
one from both dimensions and preserves the stated properties.
Inserting a zero block gives the identity correspondence.
\end{proof}

\subsection{Canonical roots and the relative coefficient}

For families on $E\times T$, write
$D(U,V)=R\operatorname{pr}_{T,*}R\mathcal Hom_E(U,V)$.
Its degrees are zero and one. Put
\[
 \lambda_n=\det D(\mathcal U_n,\mathcal U_n)=K_{\mathcal M_n},\qquad
 L_n=K_{\overline{\mathcal M}_n},\qquad W=\operatorname{Lie}(E),
\]
where $\mathcal U_n$ denotes the universal sheaf. Determinants
carry their Euler parity.

\begin{proposition}\label{gh:roots}
On every extension $0\to T_a\to T\to T_b\to0$ there is an
$E$-equivariant isomorphism
\[
 \lambda_{a+b}\simeq
 \lambda_a\otimes\lambda_b\otimes
                  \bigl(\det D(T_a,T_b)\bigr)^{\otimes2}.
\]
It squares to the threefold determinant comparison and is compatible
with all flag refinements and units. If
$a_n:\mathcal M_n\to\overline{\mathcal M}_n$ is the quotient map, then
\[
 a_n^*L_n\simeq\lambda_n\otimes\det W.
\]
The relative root $L_n\otimes\pi_n^*K_{B_E}^{-1}$ has even parity and
has its evaluation trivialization at $n=0$.
\end{proposition}
\begin{proof}
The curve tangent complex is $D(T,T)[1]$, whose cotangent determinant
is $\det D(T,T)$. The unit and the fixed trace give the two graded
pieces $\mathbb C$ and $\mathbb C[-2]$ of the Postnikov filtration of
$R\operatorname{Hom}_S(G,G)$. Determinant multiplicativity therefore
identifies the threefold endomorphism determinant with $\lambda_n^2$.
This uses a filtration of complexes, not a multiplicative splitting.

Filter both arguments of $D(T,T)$ by the extension. Determinant
multiplicativity gives the four ordered Hom blocks. Curve Serre
duality identifies
\[
 D(T_b,T_a)\simeq D(T_a,T_b)^\vee[-1].
\]
Duality and the odd shift each invert the determinant, so the two
cross determinants are identified. This gives the root isomorphism,
and its cross square is the threefold cross determinant.
For a longer flag, every bracketing refines to the determinant map of
the same doubly filtered complex. Each Hom block has Euler rank zero.
Thus the regroupings introduce no superline signs, and the root maps
themselves agree. A zero block has its distinguished determinant unit.
Translation preserves these filtrations and the chosen Serre trace.

The cotangent triangle of $a_n$ has relative cotangent line $W^*$.
It gives $a_n^*L_n=\lambda_n\otimes\det W$.
More generally, the determinant of an $r$-fold product over $B_E$
contains the inverse of $r-1$ base determinants: the correction is
$(\det W)^{1-r}$. Removing the single base root from $L_n$ leaves
relative dimension zero. At the zero object this is
$K_{B_E}\otimes K_{B_E}^{-1}$.
\end{proof}

A smooth stack $Y$ has zero $d$-critical structure: its sheaf of
critical sections is zero, as one sees on the identity embedding of a
smooth atlas. Its critical canonical line is $K_Y^{\otimes2}$.
We use the dimension-signed monodromy convention in which the
canonically oriented smooth coefficient on a stack of dimension $m$ is
\[
 \boldsymbol\Phi_Y=(\mathbb Q_Y[m],(-1)^m\operatorname{id}).
\]
The sign is part of the convention; shifting a complex alone does not
change its specified automorphism. Since both
$\overline{\mathcal M}_n$ and $B_E$ have dimension $-1$, the
\emph{normalized relative coefficient} is
\begin{equation}\label{gh:relative-coefficient}
 \boldsymbol K_n=
 \boldsymbol\Phi_{\overline{\mathcal M}_n}
       \otimes\pi_n^*\boldsymbol\Phi_{B_E}^{-1}
 =(\mathbb Q_{\overline{\mathcal M}_n},+\operatorname{id}).
\end{equation}
Its underlying complex is denoted $K_n$.
The root comparison in Proposition~\ref{gh:roots}, with the base line
removed, identifies $a_n^*\boldsymbol K_n$ with the independently
normalized coefficient
$(\mathbb Q_{\mathcal M_n},+\operatorname{id})$.
In particular $K_0=\mathbb Q_{B_E}$.
These are specified coefficients on the smooth spherical sector.
They do not identify reduced critical coefficients on a sheaf locus
with additional K3 deformation directions.

\subsection{The product, its trace, and direct sum}

All tensors and pullbacks below are derived. We use rational Betti
complexes on analytic stacks, defined by smooth descent, in an
unbounded ambient derived category. The coefficients $K_n$ are
bounded and constructible; their ordinary direct images need not be
bounded above. Put
\[
 Z_{a,b}=\overline{\mathcal M}_a\times_{B_E}
                          \overline{\mathcal M}_b,\qquad
 \mathcal A_n=R\pi_{n,*}K_n.
\]
Let $\overline{\mathcal E}_{a,b}$ be the extension stack on a common
torsor, with maps
$Z_{a,b}\xleftarrow{p}\overline{\mathcal E}_{a,b}
\xrightarrow{q}\overline{\mathcal M}_{a+b}$.
Here the first endpoint is the subobject.

\begin{proposition}\label{gh:coefficient-maps}
The complex orientation gives an isomorphism
\[
 u_{a,b}:p^*(K_a\boxtimes_{B_E}K_b)
                  \xrightarrow{\ \sim\ }q^!K_{a+b}
\]
which is the identity on the constant coefficient. Its shift is zero.
These maps satisfy every flag and unit equation, including on nonsplit
extensions, and intertwine the specified $+1$ automorphisms.
\end{proposition}
\begin{proof}
For a smooth complex stack $Y$ of dimension $m$,
$\omega_Y=\mathbb Q_Y[2m]$.
If $f:Y\to Z$ is any map between smooth stacks, transitivity of
exceptional pullback gives
$f^!\mathbb Q_Z=\mathbb Q_Y[2(\dim Y-\dim Z)]$.
The source and target of $q$ both have dimension $-1$ by
Proposition~\ref{gh:sector}. Thus $q^!K_{a+b}$ is the constant
coefficient without a shift, as is the endpoint pullback. Define
$u_{a,b}$ by the identity under these dualizing identifications.

Both substitutions for a three-step flag are made on the same full
flag stack. Their smooth endpoint legs have relative dimension zero.
Smooth base change and transitivity identify both recipients with the
constant coefficient there; both composites are its identity.
For every longer flag the same argument uses its full endpoint and
middle-object maps. The functorial adjunction and base-change
identifications give the refinement coherences. Repeated or zero
steps are identity correspondences, so both unit maps are identities.
\end{proof}

The relative external-product map will be denoted
\[
 \kappa_{a,b}:\mathcal A_a\otimes_{\mathbb Q_{B_E}}^{\mathbb L}
                  \mathcal A_b
       \longrightarrow R\pi_{Z_{a,b},*}\mathbb Q_{Z_{a,b}}.
\]
The Hall product uses this arrow in its forward direction:
\begin{equation}\label{gh:hall-product}
 \begin{split}
 \mathcal A_a\otimes^{\mathbb L}_{\mathbb Q_{B_E}}\mathcal A_b
 &\xrightarrow{\kappa_{a,b}}
                  R\pi_{Z_{a,b},*}\mathbb Q\\
 &\xrightarrow{\text{unit of }p^*\dashv Rp_*}
                  R\pi_{Z_{a,b},*}Rp_*p^*\mathbb Q\\
 &\xrightarrow{u_{a,b}}
                  R\pi_{a+b,*}Rq_*q^!K_{a+b}
 \xrightarrow{\operatorname{tr}_q}\mathcal A_{a+b}.
 \end{split}
\end{equation}
The trace is the proper counit: $q$ is representable projective, so
$Rq_*=Rq_!$. There is no integration over the nonproper sheaf stack.
The flag equations make
$\mathcal A=\bigoplus_{n\geq0}\mathcal A_n$ an associative unital
algebra over $B_E$, with unit $\mathbb Q_{B_E}$. The sum has finite
charge support, and no completion is used.

We recall the precise cohomology comparison for torsion sheaves on a
curve~\cite[\S\S3.1--3.3]{JindalLim2026}. Let
\[
 X_1=\mathcal M_1=E\times B\mathbb G_m,\qquad
 N=H^*(E;\mathbb Q)[z],\quad |z|=2,\qquad
 T_n=N^{\otimes n},\quad
 \mathsf B_n=T_n^{S_n}.
\]
The action of $S_n$ simultaneously permutes the curve factors and
Chern variables, with the Koszul signs on curve cohomology.
Ordered direct sum induces the isomorphism
\begin{equation}\label{gh:curve-cohomology}
 H^*(\mathcal M_n;\mathbb Q)
    \xrightarrow{\ \sim\ }
 \bigl(H^*(E^n;\mathbb Q)[z_1,\ldots,z_n]\bigr)^{S_n}
 =\mathsf B_n.
\end{equation}
The map matters here. The full flag stack has an associated-graded
map to $X_1^n$ whose pullback is a vector bundle stack cohomology
equivalence. The proper map to $\mathcal M_n$ identifies its
cohomology with the invariant subspace, by the small-map and Galois
comparison in the cited theorem. The split-flag section is the inverse
of the associated-graded pullback on cohomology, and its composite
with the proper map is ordered direct sum. Thus
\eqref{gh:curve-cohomology} is the pullback of that actual morphism.
We use the cited small-map comparison as a theorem, rather than
infer an isomorphism from dimensions.

\begin{proposition}[Coefficient objects and direct sum]
\label{gh:direct-sum}
The maps $\kappa_{a,b}$ are equivalences, and direct sum induces
equivalences of coefficient objects
\[
 \Phi_n:\mathcal A_n
 \xrightarrow{\ \sim\ }
 \left(\mathcal A_1^{\otimes^{\mathbb L}_{\mathbb Q_{B_E}}n}
                                                    \right)^{S_n}.
\]
For $s_{a,b}:Z_{a,b}\to\overline{\mathcal M}_{a+b}$, the actual maps
\begin{equation}\label{gh:direct-sum-map}
 d_{a,b}:\mathcal A_{a+b}
 \xrightarrow{s_{a,b}^*}
 R\pi_{Z_{a,b},*}\mathbb Q
 \xrightarrow{\kappa_{a,b}^{-1}}
 \mathcal A_a\otimes_{\mathbb Q_{B_E}}^{\mathbb L}\mathcal A_b
\end{equation}
are coherently coassociative and cocommutative for the usual Koszul
symmetry. Both zero-endpoint components are the unit identifications.
After fibre realization and graded cohomology, $d_{a,b}$ is the
inclusion
$\mathsf B_{a+b}\hookrightarrow\mathsf B_a\otimes\mathsf B_b$.
\end{proposition}
\begin{proof}
The equivariant direct-sum map
$X_1^n/S_n\to\mathcal M_n$ gives an actual pullback of constant
coefficients over $B_E$. Locally trivializing the universal torsor
identifies its stalk with the cochain map in
\eqref{gh:curve-cohomology}, so its cone has acyclic stalks.
This uses local triviality of the associated fibration, not proper
base change for $\pi_n$.

The stalk of $\kappa_{a,b}$ is the ordinary rational K\"unneth map.
The formula \eqref{gh:curve-cohomology} shows that all cohomology is
bounded below and finite-dimensional in each degree. K\"unneth
therefore applies degree by degree, including every positive Chern
power. The same cone argument proves that $\kappa_{a,b}$ is an
equivalence. Taking invariants under a finite group over $\mathbb Q$
is exact, by averaging; this proves the formula for $\Phi_n$.

Associativity and commutativity of direct sum, functorial pullback,
and the coherent external products give the assertions about $d$.
At a zero endpoint the stack is
$B_E\times_{B_E}\overline{\mathcal M}_n
=\overline{\mathcal M}_n$ and the direct-sum map is the identity.
Finally, compose with all ordered charge-one direct-sum maps. Both
routes are the same ordered direct sum, so the cohomology map is
restriction from $S_{a+b}$-invariants to
$S_a\times S_b$-invariants.
\end{proof}

\begin{proposition}[The constant classes]\label{gh:constants}
Let $\mathbf1_n$ be the constant class in
$H^0(\mathcal M_n;\mathbb Q)$. Then
\[
 \mathbf1_a*\mathbf1_b=\binom{a+b}{a}\mathbf1_{a+b},
 \qquad \mathbf1_1^{*n}=n!\mathbf1_n .
\]
\end{proposition}
\begin{proof}
Over a sheaf supported at $a+b$ distinct points, the proper extension
map has $\binom{a+b}{a}$ sheets, one for each choice of the subobject
support. The coefficient map is the identity, so its complex-oriented
trace has this degree. The stack $\mathcal M_{a+b}$ is connected:
deform each local Jordan block by a companion polynomial to distinct
nearby roots, then use the connected distinct-support configuration
space. A degree-zero class is determined by its value there.
Iteration gives the factorial.
\end{proof}

\subsection{One formality comparison for all the operations}

To compare the relative product and coproduct, it is necessary to
retain the coherent translation action. Purity provides a simultaneous
comparison of this entire diagram. We give the coefficient lift and
the extension of the purity argument that accommodates all Chern
powers.

Use rational graded-polarizable mixed Hodge structures, and write
$\mathcal C=D^b(\operatorname{MHS}_{\mathbb Q})$ for its enhanced
symmetric monoidal derived category. The mixed Hodge category of a
point is $\operatorname{Ind}(\mathcal C)$. On complex Artin stacks
we use the six-operation theory of
\cite[Theorems 3.1, 3.8 and Proposition 3.16]{Tubach2025}.
It supplies smooth descent, $\mathbb A^1$-invariance, purity,
projection formulas, and comparison with Betti realization in the
bounded-below cohomologically constructible range used here.

\begin{lemma}\label{gh:hodge-lift}
The coefficient maps, Hall products, full flag coherences, direct-sum
pullbacks and translation-action diagrams above lift together to
mixed Hodge coefficients. Their Betti realizations are the specified
maps, including their proper traces and unit normalizations.
\end{lemma}
\begin{proof}
Set $K_n^{\mathrm H}=\mathbb Q^{\mathrm H}_{\overline{\mathcal M}_n}$
with automorphism $+1$, including $K_0^{\mathrm H}=\mathbb Q^{\mathrm H}_{B_E}$.
This lifts the normalized relative coefficient
\eqref{gh:relative-coefficient}; no absolute vanishing-cycle Hodge
module or half-Tate twist is needed.
For a smooth stack $Y$ of dimension $m$,
$\omega_Y^{\mathrm H}=\mathbb Q_Y^{\mathrm H}(m)[2m]$.
Transitivity gives, for any map between smooth stacks,
\[
 f^!\mathbb Q_Z^{\mathrm H}
 =\mathbb Q_Y^{\mathrm H}(\dim Y-\dim Z)
                                      [2(\dim Y-\dim Z)].
\]
Thus the extension map $q$ has both shift and Tate twist zero.
Define $u_{a,b}^{\mathrm H}$ by the identity under this
identification. Hodge purity realizes the same complex orientation,
so its Betti map is $u_{a,b}$ with scalar exactly one.
All flag substitutions are the same identities under smooth base
change and transitivity. Enhanced functoriality supplies their
coherences.

The unit $\mathbb Q^{\mathrm H}\to Rp_*\mathbb Q^{\mathrm H}$
is an equivalence. Locally write the extension stack as $[V^1/V^0]$.
The \v Cech terms are $V^1$ times powers of $V^0$, hence vector
bundles over the endpoint base. Their units are equivalences by
$\mathbb A^1$-invariance and are compatible with all face and
degeneracy maps. Smooth descent gives the assertion. This argument
uses algebraic maps, rather than a real contraction asserted to be a
Hodge morphism.

The map $q$ is representable projective. Its Hodge counit
$Rq_!q^!\to\operatorname{id}$ realizes the actual proper trace.
Set $\mathcal A_n^{\mathrm H}=R\pi_{n,*}K_n^{\mathrm H}$.
The Hodge external-product map is an equivalence: after the smooth
atlas $\mathrm{pt}\to B_E$, its Betti realization is the K\"unneth
equivalence already proved. Betti realization is conservative on
the bounded-below cohomologically constructible objects in question.
Using the forward external-product arrow, the unit of $p$, $u^{\mathrm H}$
and the proper counit gives exactly \eqref{gh:hall-product}.
Direct sum and translation are algebraic maps, so their pullbacks
and all their composition diagrams lift in the same formalism.
\end{proof}

\begin{lemma}[Pure diagrams with unbounded Chern degree]
\label{gh:ind-formality}
Let $\mathcal P\subset\mathcal C$ be the full subcategory of bounded
objects $C$ for which $H^j(C)$ is pure of weight $j$.
The rational realization functor on $\operatorname{Ind}(\mathcal P)$
admits a monoidal natural equivalence with
\[
 h(C)=\bigoplus_j H^j(C;\mathbb Q)[-j]
\]
with zero differential. Every bounded-below object of
$\operatorname{Ind}(\mathcal C)$ with finite-dimensional cohomology
in each degree and weight $j$ in degree $j$ lies in this subcategory.
The equivalence applies to all morphisms and coherent diagrams
between these objects.
\end{lemma}
\begin{proof}
For bounded objects this is the rational pure-weight theorem of
\cite[Proposition 2.7, Lemma 4.4 and Corollary 4.6]{CiriciHorel2022}.
Its mechanism is a natural weak splitting of the weight filtration
followed by truncation in each weight. In cohomological notation, on
the weight-$w$ complex retain the terms below degree $w$ and the
cycles in degree $w$. Inclusion into the original complex and
projection onto its degree-$w$ cohomology are quasi-isomorphisms
when cohomology has weight equal to degree. Tensor products respect
this truncation: the boundary degree is the sum of the two boundary
degrees, and the tensor of cycles is a cycle. The resulting
transformations are monoidal. The rational weak splitting need not
be canonical; it is a splitting of the underlying weight-filtered
vector spaces, not a splitting of mixed Hodge structures.
Localizing these natural equivalences gives the corresponding
enhanced statement, as in the proof of
\cite[Theorem 7.1]{CiriciHorel2022}.
A quasi-isomorphism with a pure endpoint has a pure source, so this
is the full pure derived subcategory.

The category $\mathcal P$ contains the tensor unit and is closed
under tensor products by K\"unneth. Its objects are compact in
$\operatorname{Ind}(\mathcal C)$. For filtered diagrams
$C_i,D_j$ in $\mathcal P$, the mapping spaces in both Ind categories
are
\[
 \operatorname{Map}\left(\mathop{\operatorname{colim}}_i C_i,
                         \mathop{\operatorname{colim}}_j D_j\right)
 =\lim_i\mathop{\operatorname{colim}}_j\operatorname{Map}(C_i,D_j).
\]
Thus $\operatorname{Ind}(\mathcal P)\to
\operatorname{Ind}(\mathcal C)$ is fully faithful, including on
higher homotopies. Tensor preserves filtered colimits separately.
The monoidal equivalence therefore extends to this full Ind
subcategory. Exactness of filtered colimits identifies its extended
cohomology functor with $h$.

If $C$ satisfies the final hypotheses, each good truncation
$\tau_{\leq m}C$ is a bounded object with finite mixed Hodge
cohomology, hence belongs to $\mathcal C$ and to $\mathcal P$.
The canonical map
$\operatorname{colim}_m\tau_{\leq m}C\to C$ is an equivalence:
the t-structure is compatible with filtered colimits and is
separated on these bounded-below objects. No individual truncation
has been made an algebra. The monoidal equivalence is on the full
Ind subcategory, so it retains the entire diagram and its coherences.
\end{proof}

\begin{theorem}[Simultaneous Hall, direct-sum and action comparison]
\label{gh:simultaneous}
There is a rational monoidal formality comparison for the normalized
smooth spherical sector that simultaneously retains the actual Hall
products, proper flag traces, direct-sum pullbacks and translation
actions. After descent to $B_E$, its strict action model is
\[
 \Lambda_E=\bigwedge(\xi_a,\xi_b),\quad |\xi_a|=|\xi_b|=-1,
 \qquad
 \mathsf B_n=
 \left(\bigwedge(a_1,b_1,\ldots,a_n,b_n)
                   [z_1,\ldots,z_n]\right)^{S_n},
\]
with zero internal differential, primitive coproduct on
$\Lambda_E$, and actions
$I_a=\sum_i\partial_{a_i}$ and $I_b=\sum_i\partial_{b_i}$.
The degrees are $|a_i|=|b_i|=1$ and $|z_i|=2$.
The direct-sum map is the inclusion of Proposition~\ref{gh:direct-sum}.
\end{theorem}
\begin{proof}
Put $C_n^{\mathrm H}=R\Gamma(\mathcal M_n,\mathbb Q^{\mathrm H})$.
The pullback in \eqref{gh:curve-cohomology} is algebraic, hence is an
isomorphism of mixed Hodge structures. Curve cohomology is pure of
its cohomological weight and each $z_i$ has degree and weight two.
Finite rational invariants preserve purity. Thus $C_n^{\mathrm H}$
is bounded below, degreewise finite, and pure of weight $j$ in degree
$j$. The same holds for full flag objects by their associated-graded
vector bundle stack comparison, and for their products with $E^k$.

Lemma~\ref{gh:hodge-lift} places all the Hall and proper-trace maps
in this diagram. Include the actual direct-sum pullbacks and their
associativity diagrams; their sources and targets satisfy the same
purity hypotheses. Lemma~\ref{gh:ind-formality} applies to all of
them at once. In particular it compares multiplication with the
actual cohomological Hall product of
\cite[\S3.3]{JindalLim2026}. All endpoint orderings are subobject
first.

The same pure diagram contains the cup product, addition and identity
of $E$, and the translation coactions
\[
 C_n^{\mathrm H}\longrightarrow
 R\Gamma(E,\mathbb Q^{\mathrm H})\otimes C_n^{\mathrm H}.
\]
Consequently the same monoidal equivalence formalizes this Hopf
coaction and its compatibility with every Hall and direct-sum map.
The cochains of $E$ are perfect; their Hopf dual is the rational
chain action algebra. Choosing dual loop and cohomology bases gives
the exterior model $\Lambda_E$ and the contractions stated above.
For a torus this model can also be obtained by choosing primitive
based loop cycles in normalized chains: the shuffle Pontryagin
product extends them to the exterior Hopf model, and they generate
its homology. The rational Borel comparison is natural in the
action~\cite{Franz2003}. No dual of the
infinite-dimensional $B\mathbb G_m$ cochain object is taken.

Locally constant derived coefficient systems over $BE$ are
homotopy-coherent $E$-actions on their fibres. The associated
fibration defining $\pi_n$ realizes precisely that action system.
Thus the simultaneous comparison descends to the actual relative
maps. It does not use a general assertion that fibre cohomology
determines morphisms over $BE$. Composing direct sum with ordered
charge-one sums gives the stated invariant inclusion, as in
Proposition~\ref{gh:direct-sum}.
\end{proof}

The common-base cochain model of this action is
\begin{equation}\label{gh:koszul-model}
 \mathfrak t(\mathsf B_n)
 =\left(\mathbb Q[u,v]\otimes\mathsf B_n,\,
                         d=uI_a+vI_b\right),
 \qquad |u|=|v|=2.
\end{equation}
On ordered factors it has $da_i=u$ and $db_i=v$ for every $i$.
Tensoring these models over the same $\mathbb Q[u,v]$ gives the
diagonal action on the tensor product. The model is free with a
filtration by exterior degree whose differential lowers that degree,
so this tensor computes the derived tensor in the cases at hand.
The charge-zero model is $\mathbb Q[u,v]$, the cochains of $BE$.
Formality was applied before this Borel construction; it has not
removed its differential. Indeed $u,v$ have Hodge weight one and
cohomological degree two. No pure-weight semisimplicity assertion on
$BE$ is used; that assertion fails for this nonaffine stabilizer
\cite[Remark 3.25]{Tubach2025}.

\subsection{The relative diagonal and the constant-target obstruction}

Choose $a,b\in H^1(E;\mathbb Q)$ with $\int_E ab=1$.
On two ordered factors put
\begin{equation}\label{gh:diagonal}
 \delta_{12}=(a_1-a_2)(b_1-b_2)
 =a_1b_1+a_2b_2-a_1b_2+b_1a_2 .
\end{equation}
This is the fibre expression for the elliptic diagonal class
$\Delta_E$. Over the common base it is closed, since both differences
are closed in \eqref{gh:koszul-model}.
It is nonzero even there: the chain map setting
$u=v=a_1=b_1=0$ sends it to the nonzero class $a_2b_2$.
Geometrically it is the point class pulled back by relative position.
Indeed
\[
 \overline{\mathcal M}_1=B\mathbb G_m,\qquad
 \overline{\mathcal M}_1\times_{B_E}\overline{\mathcal M}_1
          \simeq E\times B\mathbb G_m^2.
\]
The extra $E$ is retained by the common torsor.

\begin{proposition}\label{gh:no-relative-rho}
There is no degree-zero, charge-preserving morphism
$\mathcal A\to\mathbb Q_{B_E}[x]$, with $x$ of cohomological degree
zero and charge one, whose fibre cohomology sends $\mathbf1_1$ to
$x$.
\end{proposition}
\begin{proof}
The map $B\mathbb G_m=\overline{\mathcal M}_1\to BE$ classifies
the trivial $E$-torsor. Thus $u\mathbf1_1=0$ in its global
hypercohomology. A morphism over $\mathbb Q_{B_E}$ induces an
$H^*(BE)$-linear map on global hypercohomology. The required
degree-zero normalization sends $\mathbf1_1$ to $x$, and would
therefore send zero to the nonzero class $ux$.

The same obstruction is visible in
$\mathbb Q[u,v]\otimes\bigwedge(a,b)[z]$, where $da=u$.
An $\mathbb Q[u,v]$-linear degree-zero chain map to
$\mathbb Q[u,v][x]$ sending $1$ to $x$ would have $f(a)=0$ and
$f(da)=ux\ne0$. The hypercohomology argument makes the
contradiction independent of this choice of representative.
\end{proof}

This obstruction concerns the constant relative target with the stated
normalization. A quotient after trivial-torsor base change and graded
cohomology has a different source and will exist below. A shifted
integration map or a target retaining the torsor cochains likewise
has a different type.
